\chapter{Kategorie pochodne, struktury Hodge'a, cykle zanikające i monodromia}
\label{ch:xiv-hodge-monodromia}

\section{Motywacja i kolejność konstrukcji}
W rozdziale~\ref{ch:ih-perverse-xiii} pojawiły się trzy zjawiska,
których nie da się opisać samymi grupami kohomologii:
obiekt $Rf_*K$ nie jest na ogół określony przez swoje snopy
kohomologii $\mathcal H^q(Rf_*K)$, kompleks przecięcia jest definiowany przez
warunki na kiełkach i kopiełkach, a w degenerującej rodzinie
rozmaitości istotne są nie tylko grupy $H^q$, ale także sposób,
w jaki zmieniają się one przy okrążaniu punktu krytycznego.

Najpierw rozwiniemy minimum rachunku w kategoriach pochodnych,
następnie strukturę Hodge'a w przypadku gładkim i jej wersję
mieszaną dla prostego przykładu osobliwego. Na koniec
zdefiniujemy cykle pobliskie oraz zanikające i obliczymy
monodromię dla osobliwości kwadratowych.
Wszystkie współczynniki, jeżeli nie zaznaczono inaczej, należą
do $\Q$; przy lokalnych obliczeniach analitycznych jawnie
używamy też $\C$ lub $\Z$.

\section{Kompleksy, odwzorowania stożkowe i trójkąty wyróżnione}
\begin{definicja}[Kompleks kohomologiczny]
Niech $\mathcal A$ będzie kategorią abelową. Kompleks
kohomologiczny $A^\bullet$ jest rodziną obiektów
$A^i\in\mathcal A$ i morfizmów
$d_A^i:A^i\to A^{i+1}$ spełniających
$d_A^{i+1}d_A^i=0$. Dla $m\in\Z$ definiujemy
przesunięcie $A[m]$ przez
$A[m]^i=A^{i+m}$ oraz $d_{A[m]}^i=(-1)^m d_A^{i+m}$.
Kohomologię oznaczamy
$H^i(A^\bullet)=\ker d_A^i/\operatorname{im}d_A^{i-1}$.
Morfizm kompleksów $f:A^\bullet\to B^\bullet$ jest
\emph{quasi-izomorfizmem}, gdy $H^i(f)$ jest izomorfizmem
dla każdego $i$.
\end{definicja}

\begin{definicja}[Stożek morfizmu]
Dla morfizmu kompleksów $f:A^\bullet\to B^\bullet$
definiujemy $\operatorname{Cone}(f)^i=B^i\oplus A^{i+1}$ oraz
\[
d_{\operatorname{Cone}(f)}(b,a)=(d_Bb+f(a),-d_Aa).
\]
Istnieje naturalny ciąg morfizmów kompleksów
$B^\bullet\to\operatorname{Cone}(f)\to A^\bullet[1]$,
a ciąg $A^\bullet\xrightarrow f B^\bullet\to
\operatorname{Cone}(f)\to A^\bullet[1]$
nazywamy standardowym trójkątem stożkowym.
\end{definicja}
\begin{lemat}[Stożek jest kompleksem]
\label{lem:cone-complex-xiv}
Dla każdego morfizmu kompleksów $f:A^\bullet\to B^\bullet$
operator $d_{\operatorname{Cone}(f)}$ spełnia $d^2=0$.
\end{lemat}
\begin{proof}
Dla $(b,a)\in B^i\oplus A^{i+1}$ otrzymujemy
$d^2(b,a)=(d_B^2b+d_Bf(a)-f(d_Aa),d_A^2a)$.
Pierwszy i ostatni składnik znikają, a środkowe znoszą się,
gdyż $f$ jest morfizmem kompleksów, tj. $d_Bf=fd_A$.
\end{proof}

\begin{twierdzenie}[Długi ciąg stożka]
\label{thm:cone-les-xiv}
Dla $f:A^\bullet\to B^\bullet$ istnieje naturalny długi
ciąg dokładny
\[
\cdots\to H^i(A)\xrightarrow{H^i(f)}H^i(B)
\to H^i(\operatorname{Cone}(f))\to H^{i+1}(A)\to\cdots.
\]
W szczególności $f$ jest quasi-izomorfizmem wtedy i tylko wtedy,
gdy $\operatorname{Cone}(f)$ jest acykliczny.
\end{twierdzenie}
\begin{proof}
Rzut na drugi czynnik daje krótki ciąg dokładny kompleksów
$0\to B^\bullet\to\operatorname{Cone}(f)\to A^\bullet[1]\to0$.
Z konstrukcji homomorfizmu łączącego otrzymujemy
$H^{i+1}(A)\to H^{i+1}(B)$ równe, z przyjętą konwencją
znaku, $H^{i+1}(f)$: istotnie reprezentant $(0,a)$
klasy z $H^i(A[1])$ ma różniczkę $(f(a),0)$.
Długi ciąg krótkiego ciągu kompleksów jest dokładnie
ciągiem z tezy. Jeśli $H^i(f)$ są izomorfizmami,
dokładność wymusza $H^i(\operatorname{Cone}(f))=0$.
Odwrotnie, gdy kohomologie stożka znikają, te same
fragmenty ciągu wymuszają jednocześnie różnowartościowość
i surjektywność wszystkich $H^i(f)$.
\end{proof}

\begin{przyklad}[Niebanalne zjawisko zapisane przez stożek]
Dla morfizmu kompleksów $\Z\xrightarrow{\cdot m}\Z$
(skupionych w stopniu $0$) stożek ma w stopniach
$-1,0$ grupy $\Z,\Z$ i różniczkę będącą mnożeniem
przez $m$. Dla $m\neq0$ jego jedyną niezerową kohomologią
jest $H^0\cong\Z/m\Z$. W ten sposób obiekt ilorazowy
można odtworzyć z morfizmu, nie wybierając rozszczepienia.
\end{przyklad}

\begin{definicja}[Kategoria pochodna]
Kategorię homotopii $K(\mathcal A)$ otrzymujemy z kategorii
kompleksów w $\mathcal A$ przez utożsamienie morfizmów
łańcuchowo homotopijnych. Jej \emph{lokalizację względem
quasi-izomorfizmów} oznaczamy $D(\mathcal A)$ i nazywamy
kategorią pochodną. $D^b(\mathcal A)$ oznacza pełną
podkategorię obiektów o ograniczonej kohomologii.
W tej kategorii obrazy trójkątów stożkowych są
\emph{trójkątami wyróżnionymi}.
\end{definicja}
\begin{uwaga}[Co pozostaje poza rachunkiem elementarnym?]
Sama definicja lokalizacji nie zapewnia automatycznie
istnienia zbiorów morfizmów o właściwej wielkości ani
formalizmu funktorów pochodnych. Dla rozważanych dalej
kategorii snopów stosujemy standardowe twierdzenia
o rezolucjach iniektywnych i kompleksach $K$-iniektywnych.
Nie utożsamiamy również każdego trójkąta wyróżnionego
z dosłownym krótkim ciągiem snopów.
\end{uwaga}

\section{Funktory pochodne i dlaczego nie wystarczają snopy \texorpdfstring{$R^qf_*$}{Rq f*}}
\begin{definicja}[Pochodny obraz prosty i hiperkohomologia]
Niech $f:X\to Y$ będzie ciągłym odwzorowaniem
lokalnie zwartych przestrzeni, a $\mathcal F^\bullet$
kompleksem snopów $k$-przestrzeni wektorowych na $X$.
Wybieramy quasi-izomorfizm
$\mathcal F^\bullet\to I^\bullet$ do kompleksu
$K$-iniektywnego. Definiujemy
$Rf_*\mathcal F^\bullet=f_*I^\bullet$ w kategorii pochodnej
oraz $R\Gamma(X,\mathcal F^\bullet)=\Gamma(X,I^\bullet)$.
Dla pojedynczego snopu wystarcza zwykła
rezolucja iniektywna.
Niech $\mathbb H^q(X,\mathcal F^\bullet)=
H^q(R\Gamma(X,\mathcal F^\bullet))$.
Dla pojedynczego snopu $\mathcal F$ mamy
$R^qf_*\mathcal F=H^q(Rf_*\mathcal F)$.
\end{definicja}
\begin{twierdzenie}[Ciąg spektralny Leraya]
\label{thm:leray-derived-xiv}
Niech $f:X\to Y$ będzie ciągłym odwzorowaniem
przestrzeni topologicznych, a $\mathcal F$ snopem grup
abelowych na $X$. Dla zwykłej kohomologii snopowej
i nieujemnych funktorów pochodnych otrzymujemy
pierwszoćwiartkowy, zbieżny do kohomologii w każdym
stopniu całkowitym ciąg spektralny
$E_2^{p,q}=H^p(Y,R^qf_*\mathcal F)
\Rightarrow H^{p+q}(X,\mathcal F)$.
\end{twierdzenie}
\begin{proof}[Źródło twierdzenia]
Kompozycja $\Gamma(Y,-)\circ f_*$ jest równa
$\Gamma(X,-)$. Twierdzenie o ciągu spektralnym
złożenia funktorów pochodnych (wraz z warunkiem
acykliczności obrazów rezolwent iniektywnych)
daje wskazaną stronę $E_2$ i jej zbieżność.
Sama konstrukcja stożka nie zastępuje tego
twierdzenia o rezolwentach.
\end{proof}
\begin{uwaga}[Dane kohomologiczne nie określają obiektu]
W kategorii pochodnej nie jest na ogół prawdą, że
$K\simeq\bigoplus_qH^q(K)[-q]$.
Taki rozkład może być blokowany przez nietrywialne
morfizmy łączące w trójkątach truncacji.
Dlatego obliczenie samych $R^qf_*\mathcal F$ nie dowodzi
rozkładu $Rf_*\mathcal F$; to właśnie istotna treść
twierdzenia z rozdziału~\ref{ch:ih-perverse-xiii}.
\end{uwaga}

\begin{przyklad}[Te same snopy kohomologii, różne obiekty pochodne]
\label{ex:derived-not-split-xiv}
Niech $X=S^2$, a $\Q_X$ będzie snopem stałym.
Standardowa identyfikacja
$\operatorname{Hom}_{D(X;\Q)}(\Q_X,\Q_X[2])
\cong H^2(S^2;\Q)\cong\Q$
pozwala wybrać niezerowy morfizm
$e:\Q_X\to\Q_X[2]$.
Zdefiniujmy $K_e=\operatorname{Cone}(e)[-2]$.
Z długiego ciągu stożka wynika, że jedynymi
niezerowymi snopami kohomologii $K_e$ są
$H^0(K_e)\cong\Q_X$ i $H^1(K_e)\cong\Q_X$.
Dla $e=0$ otrzymujemy
$K_0\simeq\Q_X\oplus\Q_X[-1]$.
Dla $e\ne0$ trójkąt
$\Q_X\to K_e\to\Q_X[-1]\xrightarrow{e[-1]}\Q_X[1]$
nie jest rozszczepiony, ponieważ jego morfizm łączący
jest niezerowy. Jest to trójkąt standardowej truncacji
$K_e$, więc $K_e$ nie jest izomorficzny z sumą
$\Q_X\oplus\Q_X[-1]$.
Widzimy dwa obiekty z izomorficznymi snopami
kohomologii, lecz różnymi klasami izomorfizmu
w kategorii pochodnej.

Różnicę można też \emph{obliczyć} w hiperkohomologii.
Niech $e$ będzie generatorem $H^2(S^2;\Q)$. Po zastosowaniu
$R\Gamma(S^2,-)$ do trójkąta
$\Q_{S^2}\to K_e\to\Q_{S^2}[-1]\to\Q_{S^2}[1]$
łączący homomorfizm
$H^1(S^2;\Q_{S^2}[-1])\to H^2(S^2;\Q_{S^2})$
jest mnożeniem przez niezerowy skalar w~$\Q$, a więc
izomorfizmem $\Q\to\Q$. Długi ciąg daje dokładnie
$\mathbb H^0(S^2,K_e)\cong\mathbb H^3(S^2,K_e)\cong\Q$
oraz zanik w pozostałych stopniach. Natomiast dla $e=0$
mamy $K_0\simeq\Q_{S^2}\oplus\Q_{S^2}[-1]$, skąd
$\mathbb H^j(S^2,K_0)\cong\Q$ dla $j=0,1,2,3$.
Zatem nawet hiperkohomologia odróżnia te kompleksy,
mimo że ich \emph{snopy kohomologii} są identyczne.
\end{przyklad}

\section{Struktury Hodge'a: od rozkładu form do filtracji}
W zapisie $V_\Q\otimes_\Q\C$ korzystamy ze zwykłego
iloczynu tensorowego przestrzeni liniowych.
Jego konstrukcja i własność uniwersalna zostały już
szczegółowo przedstawione w
\href{https://github.com/dariusz-piekarz/mathematical_scipts/blob/biblioteka-matematyczna-iteracja-6/chapters/pl/differential_geometry_and_topology/03-algebra-tensorowa-i-zewnetrzna.tex}{rozdziale 3
\emph{Geometrii różniczkowej}:
\emph{Algebra tensorowa i zewnętrzna}}.
% Oryginalne etykiety DG: ch:algebra-tensorowa,
% sec:iloczyn-tensorowy, thm:tensor-hom.
Nie będziemy powtarzać tego rachunku.
\begin{definicja}[Czysta struktura Hodge'a]
Niech $m\in\Z$. \emph{Racjonalną czystą strukturą
Hodge'a wagi $m$} na skończeniewymiarowej przestrzeni
$V_{\Q}$ nazywamy rozkład
$V_{\C}=V_{\Q}\otimes_{\Q}\C=\bigoplus_{p+q=m}V^{p,q}$,
w którym $\overline{V^{p,q}}=V^{q,p}$; kreska oznacza
sprzężenie względem struktury racjonalnej.
Filtracją Hodge'a nazywamy
$F^pV_{\C}=\bigoplus_{r\ge p}V^{r,m-r}$.
\end{definicja}
\begin{stwierdzenie}[Rozkład a dwie filtracje]
\label{prop:hodge-filtration-xiv}
Dla czystej struktury Hodge'a wagi $m$ zachodzi
$V_{\C}=F^pV_{\C}\oplus\overline{F^{m-p+1}V_{\C}}$
na każdym $p\in\Z$. Odwrotnie, malejąca skończona
filtracja $F^\bullet V_{\C}$ spełniająca te równości
określa jednoznacznie strukturę Hodge'a wagi $m$ przez
$V^{p,q}=F^pV_{\C}\cap\overline{F^qV_{\C}}$
dla $p+q=m$.
\end{stwierdzenie}
\begin{proof}
W pierwszą stronę $F^p$ składa się bezpośrednia suma
wszystkich $V^{r,m-r}$ z $r\ge p$. Z kolei
$\overline{F^{m-p+1}}$ jest sumą składników
$V^{r,m-r}$ o $r\le p-1$. Składniki są rozłączne
i razem wyczerpują wszystkie pary $r+(m-r)=m$.
W drugą stronę oznaczmy
$G_p=\overline{F^{m-p}V_{\C}}$.
Filtracja $G_\bullet$ jest rosnąca, a hipoteza
przyjmuje postać $V_{\C}=F^pV_{\C}\oplus G_{p-1}$.
Stosując ją dla $p+1$, mamy
$V_{\C}=F^{p+1}V_{\C}\oplus G_p$.
Każdy $v\in F^pV_{\C}$ rozkłada się więc
jednoznacznie jako $v=u+w$,
gdzie $u\in F^{p+1}V_{\C}$ i $w\in G_p$.
Ponieważ $F^{p+1}\subseteq F^p$, także
$w=v-u\in F^p$, a zatem
\[
F^pV_{\C}=F^{p+1}V_{\C}
\oplus(F^pV_{\C}\cap G_p).
\]
Skończoność filtracji pozwala iterować to
rozszczepienie i uzyskać
$V_{\C}=\bigoplus_p(F^pV_{\C}\cap G_p)$.
W szczególności składnik indeksu $p$ jest
$F^pV_{\C}\cap\overline{F^{m-p}V_{\C}}$.
Sprzężenie zespolone zamienia ten składnik
ze składnikiem indeksu $m-p$, co daje
dokładnie rozkład Hodge'a wagi $m$.
Jednoznaczność wynika z zapisanej postaci
każdego składnika jako przecięcia.
\end{proof}

\begin{twierdzenie}[Rozkład Hodge'a dla rozmaitości zwartej Kählerowskiej]
\label{thm:kahler-hodge-xiv}
Niech $M$ będzie zwartą rozmaitością Kählerowską
wymiaru zespolonego $d$. Dla każdego $k\ge0$
istnieje naturalny rozkład
\[
H^k(M;\C)=\bigoplus_{p+q=k}H^{p,q}(M),
\qquad H^{p,q}(M)\cong H^q(M,\Omega_M^p),
\]
gdzie $\Omega_M^p$ jest snopem holomorficznych $p$-form.
Ponadto $\overline{H^{p,q}(M)}=H^{q,p}(M)$.
Jeżeli $M$ jest gładką zespoloną odmianą projektową,
to $H^k(M;\Q)$ z tym rozkładem jest czystą
strukturą Hodge'a wagi $k$.
\end{twierdzenie}
\begin{proof}[Dowód z jawnie wskazanymi twierdzeniami analitycznymi]
Na rozmaitości zespolonej formy zespolone rozkładają
się według typu: $\Omega^k(M;\C)=
\bigoplus_{p+q=k}\Omega^{p,q}(M)$.
Operator zewnętrzny ma postać
$d=\partial+\bar\partial$,
gdzie $\partial$ i $\bar\partial$ zwiększają
odpowiednio pierwszy i drugi stopień.

Z twierdzenia Hodge'a dla form różniczkowych
każda klasa de Rhama ma dokładnie jednego
reprezentanta harmonicznego dla ustalonej metryki.
Z tożsamości Kählerowskich (przyjęty tu wynik
analityczny) otrzymujemy
$\Delta_d=2\Delta_{\bar\partial}$.
Laplasjan $\Delta_{\bar\partial}$ zachowuje typ
$(p,q)$, więc rozkład form harmonicznych według
typu daje rozkład kohomologii de Rhama.
Twierdzenie Dolbeaulta identyfikuje przestrzeń
form $\bar\partial$-harmonicznych typu $(p,q)$
z $H^q(M,\Omega_M^p)$.
Sprzężenie form harmonicznych zamienia typy
$(p,q)$ i $(q,p)$. Dla rozmaitości projektowej
istnienie metryki Kählerowskiej jest klasycznym
wynikiem geometrycznym. Zatem wszystkie warunki
definicji struktury Hodge'a są spełnione.
\end{proof}

\begin{przyklad}[Prosta rzutowa i krzywa eliptyczna]
Na $\mathbb P^1$ niezerowe grupy Hodge'a to
$H^{0,0}(\mathbb P^1)\cong\C$ oraz
$H^{1,1}(\mathbb P^1)\cong\C$; w szczególności
$H^1(\mathbb P^1;\C)=0$.
Jeżeli $E=\C/\Lambda$ jest krzywą eliptyczną,
to $H^1(E;\C)=H^{1,0}(E)\oplus H^{0,1}(E)$,
oba składniki są jednowymiarowe: reprezentują je
odpowiednio klasy form $\dd z$ i $\dd\bar z$.
W stopniu drugim mamy $H^{1,1}(E)\cong\C$.
Stąd $h^{p,q}=h^{q,p}$, a $b_1(E)=2$.
\end{przyklad}

\begin{wniosek}[Parzystość nieparzystych liczb Bettiego]
Jeżeli $M$ jest zwartą rozmaitością Kählerowską,
to $b_{2r+1}(M)$ jest parzyste dla każdego $r\ge0$.
\end{wniosek}
\begin{proof}
Dla $p+q=2r+1$ nie zachodzi $p=q$, więc
składniki $H^{p,q}$ i $H^{q,p}$ tworzą rozłączne
pary tej samej liczby wymiarów. Sumowanie po
wszystkich parach daje parzysty wymiar $H^{2r+1}(M;\C)$.
\end{proof}

\section{Ciąg spektralny Hodge'a--de Rhama i jego zakres}
Niech $X$ będzie gładką zespoloną odmianą algebraiczną.
Filtracja ``głupia'' $F^p\Omega_X^\bullet=
\Omega_X^{\ge p}$ prowadzi do ciągu spektralnego
\[
E_1^{p,q}=H^q(X,\Omega_X^p)
\Longrightarrow\mathbb H^{p+q}(X,\Omega_X^\bullet).
\]
W przypadku gładkim projektowym jest to algebraiczna
wersja ciągu Hodge'a--de Rhama, a jego zbieżna
hiperkohomologia odpowiada kohomologii singularnej
$H^{p+q}(X^{\mathrm{an}};\C)$ po wykorzystaniu
twierdzenia porównawczego Grothendiecka.
\begin{twierdzenie}[Degeneracja na stronie $E_1$]
Niech $X$ będzie gładką właściwą odmianą algebraiczną
nad $\C$. Wtedy ciąg spektralny Hodge'a--de Rhama
degeneruje na stronie $E_1$.
\end{twierdzenie}
\begin{proof}[Zakres]
Dla odmian projektowych nad $\C$ jest to konsekwencja
rozkładu Hodge'a i zgodności filtracji Dolbeaulta.
W ogólnym gładkim właściwym przypadku stosuje się
odpowiednie twierdzenie analityczne lub algebraiczne.
Nie utożsamiamy tej degeneracji z dowolnym
ciągiem spektralnym dla przestrzeni osobliwych.
\end{proof}

\section{Mieszane struktury Hodge'a i osobliwa krzywa węzłowa}
\begin{definicja}[Mieszana struktura Hodge'a]
Mieszaną strukturą Hodge'a na skończeniewymiarowym
$V_{\Q}$ nazywamy rosnącą, skończoną filtrację
$W_\bullet V_{\Q}$ (filtrację wag) oraz malejącą,
skończoną filtrację $F^\bullet V_{\C}$ takie, że
na $\operatorname{Gr}_m^WV=W_mV/W_{m-1}V$
filtracja indukowana przez $F$, określona wzorem
\[
F^p(\operatorname{Gr}_m^WV)_{\C}
=\frac{(F^pV_{\C}\cap (W_mV)_{\C})+(W_{m-1}V)_{\C}}
{(W_{m-1}V)_{\C}},
\]
określa czystą strukturę Hodge'a wagi $m$
dla każdego $m$.
\end{definicja}
\begin{twierdzenie}[Deligne: zakres użycia]
Dla dowolnej zespolonej odmiany algebraicznej $X$
skończonego typu grupy $H^k(X;\Q)$ są naturalnie
wyposażone w mieszane struktury Hodge'a, funktorialne
dla morfizmów algebraicznych. Dla $X$ gładkiej
projektowej struktura na $H^k$ jest czysta wagi $k$.
\end{twierdzenie}
\begin{proof}[Status]
To głębokie twierdzenie Deligne'a. Jego dowód
wymaga zwartyfikacji, rozdzielczości osobliwości
oraz mieszanych kompleksów Hodge'a. W niniejszym
rozdziale wyprowadzamy natomiast konkretny rachunek
wag w przykładzie, przyjmując funktorialność
struktur mieszanych i ścisłość ich morfizmów.
\end{proof}

\begin{przyklad}[Waga zerowa w $H^1$ krzywej z węzłem]
\label{ex:nodal-mhs-xiv}
Niech $C$ będzie nieredukowalną rzutową krzywą
wymierną z jednym zwyczajnym węzłem, a
$\nu:\widetilde C\cong\mathbb P^1\to C$ jej normalizacją.
Preobraz węzła składa się z dwóch punktów $a,b$.
Istnieje krótki ciąg snopów stałych na $C$:
\[
0\to\Q_C\to\nu_*\Q_{\widetilde C}
\xrightarrow{\,\mathrm{ev}_a-\mathrm{ev}_b\,}
\Q_{\{x\}}\to0.
\]
Dokładność w $x$ wynika z odwzorowania $\Q\to\Q^2$
przyporządkowującej stałej $c$ parę $(c,c)$,
której kokernelem jest $\Q$ przez różnicę wartości.
Poza węzłem $\nu$ jest izomorfizmem.
W długim ciągu kohomologii $H^0(C;\Q)\to
H^0(\widetilde C;\Q)$ jest izomorfizmem,
a różnica wartości globalnej funkcji stałej
wynosi zero. Dlatego
$H^1(C;\Q)\cong\Q$,
$H^2(C;\Q)\cong\Q$.

W tym miejscu korzystamy z dokładności
w kategorii mieszanych struktur Hodge'a długiego ciągu
otrzymanego z algebraicznego kwadratu normalizacji;
jest to standardowa konsekwencja konstrukcji Deligne'a
przez zstępowanie kohomologiczne, nie wynik samej
funktorialności dla pojedynczych morfizmów.
W szczególności homomorfizm łączący
$H^0(\{x\};\Q)\to H^1(C;\Q)$ jest izomorfizmem
\emph{mieszanych struktur Hodge'a}.
Skoro jego dziedzina jest czysta wagi $0$,
również $H^1(C;\Q)$ ma wagę $0$, a nie $1$.
Tymczasem kompleks przecięcia normalizuje tę
anomalię w innym sensie:
$IH^1(C;\Q)=H^1(\mathbb P^1;\Q)=0$,
co obliczyliśmy już w rozdziale~\ref{ch:ih-perverse-xiii}.
To są dwie różne konstrukcje, a nie dwie nazwy
tego samego niezmiennika.
\end{przyklad}

\section{Rodziny, lokalne systemy i monodromia}
\begin{definicja}[Monodromia rodziny]
Niech $p:E\to B$ będzie lokalnie trywialną wiązką
z włóknem $F$, $b\in B$. Transport po pętlach
w $B$ wyznacza (po wybraniu identyfikacji włókna)
reprezentację
$\rho_q:\pi_1(B,b)\to\operatorname{Aut}(H^q(F;\Q))$,
zwaną reprezentacją monodromii kohomologicznej.
\end{definicja}
\begin{twierdzenie}[Ehresmann: wersja używana]
Jeżeli $p:E\to B$ jest surjektywną, właściwą gładką
submersją między rozmaitościami gładkimi bez brzegu,
to $p$ jest lokalnie trywialną wiązką gładką.
W szczególności $R^qp_*\Q_E$ jest lokalnym
systemem o włóknie $H^q(F;\Q)$.
\end{twierdzenie}
\begin{proof}[Szkic z precyzyjnym użyciem właściwości]
Metryka na $E$ pozwala wybrać gładki rozkład
$TE=\ker(dp)\oplus H$. Dla pola wektorowego
na małym dysku w $B$ podnosimy je poziomo do $H$.
Właściwość odwzorowania zapewnia, że nad odpowiednio
małym zwartym otoczeniem podstawy trajektorie
podniesionych pól nie uciekają w skończonym czasie.
Przepływy budują dyfeomorfizmy lokalnych włókien.
Pełny dowód uwzględnia zgodność przy zmianie pól;
przyjmujemy tutaj standardową wersję twierdzenia
Ehresmanna.
\end{proof}

\begin{center}
\begin{tikzpicture}[scale=0.85,>=Stealth,font=\small]
\draw[manifoldblue,thick] (0,0) circle (1.1);
\fill[manifoldred] (0,0) circle (2.1pt);
\node[below right] at (0,0) {$0$};
\draw[manifoldgreen,very thick,->] (1.1,0)
 arc[start angle=0,end angle=305,radius=1.1];
\node[right] at (1.25,0.75) {$\gamma$};
\node at (0,-1.42) {$\Delta^*$};
\draw[manifoldblue,thick] (4.1,0) ellipse (0.95 and 0.44);
\draw[manifoldblue,thick] (4.1,0.2) ellipse (0.95 and 0.44);
\node[above] at (4.1,1.0) {$F_t$};
\draw[->,thick] (1.5,0)--(2.8,0)
 node[midway,above] {transport};
\node at (4.1,-1.42) {$T$};
\end{tikzpicture}
\par\medskip\small Transport wzdłuż zamkniętej pętli w podstawie
wyznacza działanie monodromii na kohomologii włókna.
\end{center}

\begin{uwaga}[Punkt krytyczny]
Dla $f:X\to\Delta$, gdzie $\Delta\subset\C$ jest dyskiem,
submersja może zawieść nad $0$. Wtedy
$R^qf_*\Q$ może nie być lokalnym systemem na całym dysku,
choć jego ograniczenie do $\Delta^*=\Delta\setminus\{0\}$
może być lokalnie stałe. Monodromia pętli otaczającej
$0$ opisuje część informacji traconej przy specjalizacji.
\end{uwaga}

\section{Włókna Milnora, cykle pobliskie i zanikające}
\begin{definicja}[Włókno Milnora]
Niech $f:(\C^{n+1},0)\to(\C,0)$ będzie holomorficznym
germem, którego jedynym punktem krytycznym w
pewnym otoczeniu jest $0$, i niech $f(0)=0$.
Dla dostatecznie małych $0<|t|\ll\delta$ definiujemy
lokalne włókno Milnora
$F_t=f^{-1}(t)\cap B_\delta^{2n+2}$.
Z twierdzenia Milnora wynika, że dla odpowiednio
wybranych $t,\delta$ jego typ homotopijny nie
zależy od parametrów, a rodzina nad małym
przebitym dyskiem jest lokalnie trywialna.
\end{definicja}

\begin{definicja}[Cykle pobliskie i zanikające]
\label{def:nearby-vanishing-xiv}
Niech $f:X\to\Delta$ będzie analitycznym odwzorowaniem,
$X_0=f^{-1}(0)$, $i:X_0\hookrightarrow X$, a
$\widetilde{\Delta^*}\to\Delta^*$ uniwersalnym nakryciem.
Niech $\widetilde X^*=X\times_\Delta\widetilde{\Delta^*}$
z odwzorowaniem $\widetilde j:\widetilde X^*\to X$.
Dla $K\in D_c^b(X;\Q)$ definiujemy
\emph{kompleks cykli pobliskich}
$R\psi_f K=i^*R\widetilde j_*\widetilde j^*K$.
Istnieje naturalny morfizm specjalizacji
$\mathrm{sp}:i^*K\to R\psi_fK$.
\emph{Kompleks cykli zanikających} definiujemy jako
$R\phi_fK=\operatorname{Cone}(\mathrm{sp})$.
Wtedy istnieje trójkąt wyróżniony
\[
i^*K\longrightarrow R\psi_fK
\longrightarrow R\phi_fK\longrightarrow i^*K[1].
\]
Przy tej konwencji nie wprowadzamy dodatkowego
przesunięcia $[-1]$ do definicji $R\phi_f$.
\end{definicja}
\begin{uwaga}[Precyzja konstrukcji]
W definicji cykli pobliskich $\widetilde j$
nie jest zwykłym włożeniem otwartym:
uwzględnia uniwersalne nakrycie $\Delta^*$.
Działanie grupy jego automorfizmów
$\pi_1(\Delta^*)\cong\Z$ daje naturalny
operator monodromii $T$ na $R\psi_fK$ i $R\phi_fK$.
Otrzymanie tych funktorów w kategorii
konstruktybilnych kompleksów oraz zgodność
z wybranym nakryciem jest standardowym twierdzeniem
teorii snopów, z którego tu korzystamy.
\end{uwaga}

\begin{stwierdzenie}[Włókno Milnora i cykle zanikające]
\label{prop:milnor-stalk-xiv}
Jeśli $X$ jest małą kulą w $\C^{n+1}$,
$f:X\to\Delta$ ma izolowany punkt krytyczny $0$,
a $K=\Q_X$, to
$H^q((R\psi_fK)_0)\cong H^q(F_t;\Q)$.
Ponieważ $X_0\cap B_\delta$ jest lokalnie
ściągalne do $0$, dla $q>0$ zachodzi
$H^q((R\phi_fK)_0)\cong\widetilde H^q(F_t;\Q)$;
w stopniu $0$ analogiczne stwierdzenie dotyczy
kohomologii zredukowanej, o ile $F_t$ jest niepuste.
\end{stwierdzenie}
\begin{proof}
Pierwszy izomorfizm jest twierdzeniem o kiełkach
cykli pobliskich; używa lokalnej trywialności
fibracji Milnora i właściwych warunków
konstruktybilności. Następnie bierzemy
ciąg długi trójkąta z
Definicji~\ref{def:nearby-vanishing-xiv}.
Kiełek $i^*\Q_X$ w $0$ ma kohomologię $\Q$
w stopniu $0$ i zero w pozostałych stopniach.
Morfizm $H^0(i^*K)\to H^0(R\psi_fK)$
wysyła $1$ na funkcję stałą $1$ na każdej
składowej włókna. Dokładność daje właśnie
kohomologię zredukowaną włókna Milnora,
z uwzględnieniem przesunięć w trójkącie.
\end{proof}

\section{Dwa jawne rachunki monodromii}
\begin{przyklad}[Węzeł płaski $f(z,w)=zw$]
\label{ex:node-monodromy-xiv}
Dla $f:\C^2\to\C$, $f(z,w)=zw$, włókno
$F_t$ dla $t\ne0$ jest homotopijnie równoważne
okręgowi $S^1$: z równania $w=t/z$ wynika,
że w małej kuli $z$ przebiega pierścień.
Stąd $\widetilde H^1(F_t;\Q)\cong\Q$
i inne zredukowane grupy znikają.
Zatem jedyną niezerową kohomologią
$(R\phi_f\Q_{\C^2})_0$ jest $H^1\cong\Q$.

Monodromię można zobaczyć bezpośrednio:
przy przejściu $t\mapsto te^{i\theta}$
transport $(z,w)\mapsto(e^{i\theta}z,w)$
zachowuje równanie włókna i normę modułu $z$;
po pełnym obrocie $\theta=2\pi$ jest identycznością.
Indukowana monodromia na $H^1(F_t;\Q)$ to więc $T=1$.
Ten \emph{lokalny} wynik nie oznacza, że monodromia
na całym $H^1$ gładkiej krzywej degenerującej
do krzywej węzłowej musi być identycznością:
może ona być niebanalną transwekcją.
\end{przyklad}

\begin{twierdzenie}[Osobliwość kwadratowa $A_1$: pełne obliczenie]
\label{thm:a1-monodromy-xiv}
Niech $n\ge1$ oraz
$f:\C^{n+1}\to\C$,
$f(z_0,\dots,z_n)=\sum_{j=0}^nz_j^2$.
Wówczas lokalne włókno Milnora ma typ homotopijny $S^n$,
a monodromia geometryczna wywołuje na
$\widetilde H^n(F_t;\Z)\cong\Z$ mnożenie przez
$(-1)^{n+1}$.
\end{twierdzenie}
\begin{proof}
Po zastąpieniu $t\ne0$ przez $1$ za pomocą
skalowania zespolonego zbadamy afiniczną
kwadrykę $Q=f^{-1}(1)$. Zapiszmy $z=x+iy$,
gdzie $x,y\in\R^{n+1}$. Równanie $f(z)=1$
jest równoważne układowi
$\lVert x\rVert^2-\lVert y\rVert^2=1$,
$\langle x,y\rangle=0$.
Dla $u\in S^n$ i $p\in T_u^*S^n$ utożsamionego
z wektorem $p\perp u$, wzór
\[
(u,p)\longmapsto\sqrt{1+\lVert p\rVert^2}\,u+ip
\]
jest dyfeomorfizmem $T^*S^n\to Q$.
Odwrotność ma postać
$u=x/\lVert x\rVert$, $p=y$;
równania kwadryki zapewniają $p\perp u$.
Skalowanie $p\mapsto sp$ dla $0\le s\le1$
przekształca się w deformacyjną retrakcję
$Q\to S^n$ (przekrój zerowy).
Dla włókna lokalnego po przeskalowaniu $t$
uzyskujemy część kwadryki $Q$ ograniczoną
odpowiednio dużą kulą. Retrakcja zmniejsza
normę $\lVert z\rVert$, zatem pozostaje w tej
części i wykazuje jej ten sam typ homotopijny.

Ponieważ $f$ jest jednorodny stopnia $2$,
pętli $t(s)=e^{2\pi is}$ odpowiada transport
$z\mapsto e^{\pi is}z$, a po pełnym okrążeniu
otrzymujemy $z\mapsto-z$. Na retrakcie
$S^n\subset Q$ jest to odwzorowanie
antypodyczne stopnia $(-1)^{n+1}$.
Ponieważ retrakcja jest równoważnością
homotopijną, taki jest również operator
monodromii na $\widetilde H^n(F_t;\Z)$.
\end{proof}
\begin{przyklad}[Zmiana parzystości]
Dla $n=1$, odpowiadającego zwykłemu węzłowi
płaskiemu, monodromia działa na $H^1(S^1;\Z)$
jak $+1$. Dla $n=2$ działa na
$H^2(S^2;\Z)$ jak $-1$. Zmiana znaku wynika z
parzystości wymiaru sfery, a nie ze zmiany
liczby cykli zanikających.
\end{przyklad}

\begin{uwaga}[Różne znaczenia monodromii przy węźle]
Są tu trzy różne obiekty: monodromia lokalnego włókna Milnora,
monodromia całego $H^1$ gładkiej krzywej oraz monodromia
po zmianie parametru rodziny. Pierwsza dla $zw=t$ działa
identycznie na $H^1$ pierścienia; druga może być
nietrywialną transwekcją; po podstawieniu $t=s^m$
trzecia jest $m$-tą potęgą monodromii pierwotnej.
Nie wolno więc wnioskować z lokalnej identyczności, że
monodromia globalnej rodziny jest trywialna.
\end{uwaga}

\begin{przyklad}[Cykle zanikające i monodromia dla $f(z)=z^m$]
\label{prz:monodromy-zm-v14}
Niech $m\ge2$, a $f:\C\to\C$ będzie odwzorowaniem
$f(z)=z^m$. Jedynym punktem krytycznym w otoczeniu
zera jest $0$, ponieważ $f'(z)=mz^{m-1}$.
Dla małego $t\ne0$ włókno Milnora składa się
z dokładnie $m$ punktów
\[
 z_j(t)=|t|^{1/m}
   \exp\!\left(\frac{\mathrm i(\arg t+2\pi j)}m\right),
 \qquad j=0,\dots,m-1.
\]
Są to punkty dyskretne, więc
$H^0(F_t;\Q)\cong\Q^m$ i
$\widetilde H^0(F_t;\Q)\cong\Q^{m-1}$;
pozostałe zredukowane grupy znikają.
Zgodnie ze Stwierdzeniem~\ref{prop:milnor-stalk-xiv}
kompleks cykli zanikających ma w zerze
jedną niezerową grupę kohomologii,
w stopniu $0$ i rzędu $m-1$.
Ten przykład dotyczy wymiaru zespolonego
włókna $n=0$ i dlatego nie należy go mylić
z obliczeniem stopnia $n\ge1$ dla
osobliwości kwadratowej.

Niech $t$ okrąża dodatnio zero:
$t(s)=t_0e^{2\pi\mathrm i s}$, $0\le s\le1$.
Kontynuacja każdej gałęzi pierwiastka
prowadzi od $z_j(t_0)$ do $z_{j+1}(t_0)$
(indeksy liczymy modulo $m$).
Transport na $H_0(F_{t_0};\Q)$
jest więc cykliczną permutacją
$e_j\mapsto e_{j+1}$. Na części zredukowanej,
czyli w jądrze sumowania współczynników
$\sum_j a_j$, ma wartości własne
\[
 \zeta_m,\zeta_m^2,\dots,\zeta_m^{m-1},
 \qquad \zeta_m=e^{2\pi\mathrm i/m}.
\]
Istotnie, pełna reprezentacja permutacyjna
ma wielomian charakterystyczny $\xi^m-1$,
gdzie $\xi$ jest zmienną formalną, odróżnioną
od operatora monodromii $T$;
jej jednowymiarowa podprzestrzeń
stałych wektorów ma wartość własną $1$,
a dopełnienie sumy współczynników równej
zeru usuwa dokładnie ten jeden czynnik.
Zatem wielomian charakterystyczny na
$\widetilde H_0(F_{t_0};\Q)$ wynosi
\[
 \frac{\xi^m-1}{\xi-1}
=1+\xi+\cdots+\xi^{m-1}.
\]
Podstawienie $\xi=1$ pokazuje ponadto, że
$\det(I-T|_{\widetilde H_0(F_{t_0};\Q)})=m$.
Wersja całkowita tego rachunku wykrywa torsję:
na grupie $\widetilde H_0(F_{t_0};\Z)$ bazą są
$b_j=e_j-e_{j+1}$ dla $0\le j\le m-2$.
W ilorazie przez obraz $T-I$ wszystkie klasy $b_j$
są równe, ponieważ $(T-I)b_j=b_{j+1}-b_j$
(dla $j\le m-3$). Ostatnia relacja daje
$m[b_0]=0$, gdyż
$(T-I)b_{m-2}=(-b_0-\cdots-b_{m-2})-b_{m-2}$.
Stąd kokernelem $T-I$ jest $\Z/m\Z$;
nie pojawia się dalsza relacja, ponieważ jego rząd
wynosi $|\det(I-T)|=m$.
Tak więc $T-I$ jest izomorfizmem nad $\Q$, ale nie
nad $\Z$. To rozróżnia informację racjonalną
od całkowitej w monodromii.
Monodromia na $\widetilde H^0$ jest
odwrotnym operatorem dualnym i ma ten sam
wielomian charakterystyczny, ponieważ
zbiór pierwiastków jest niezmienniczy
względem odwracania.

Dla $m=3$, w bazie zredukowanej homologii
$b_1=e_0-e_1$, $b_2=e_1-e_2$, transport ma macierz
\[
 [T]_{(b_1,b_2)}
 =\begin{pmatrix}0&-1\\1&-1\end{pmatrix},
 \qquad T^3=I,\quad T\ne I.
\]
Jego wartości własne są pierwiastkami
pierwotnymi stopnia trzeciego.
Jest więc \emph{quasi-unipotentny},
bo $T^3=I$, lecz \emph{nie jest unipotentny}:
żadna niezerowa potęga $(T-I)$ nie znika.
To samodzielny kontrprzykład wobec
zamieniania obu pojęć bez hipotezy
o zmianie bazy parametru.
\end{przyklad}

\begin{figure}[!htbp]
\centering
\begin{tikzpicture}[>=Stealth,font=\small]
 \draw[thick,manifoldblue!70!black] (0,0) circle (1.65);
 \foreach \a/\j in {0/0,120/1,240/2} {
   \fill[manifoldred] (\a:1.65) circle (3pt);
   \node[font=\scriptsize,above right] at (\a:1.65)
        {$z_{\j}$};
 }
 \draw[->,very thick,manifoldgreen!70!black]
    (0:1.95) arc[start angle=0,end angle=105,radius=1.95];
 \draw[->,very thick,manifoldgreen!70!black]
    (120:1.95) arc[start angle=120,end angle=225,radius=1.95];
 \draw[->,very thick,manifoldgreen!70!black]
    (240:1.95) arc[start angle=240,end angle=345,radius=1.95];
 \node[right,align=left] at (2.5,.7)
      {$f(z)=z^3$};
 \node[right,align=left] at (2.5,0)
      {$z_0\mapsto z_1\mapsto z_2\mapsto z_0$};
 \node[right,align=left] at (2.5,-.7)
      {$T^3=I$, ale $T\ne I$};
\end{tikzpicture}
\caption{Monodromia trójpunktowego włókna Milnora
$f(z)=z^3$. Po dodatnim okrążeniu wartości
krytycznej punkty włókna przechodzą cyklicznie
jeden w drugi.}
\label{fig:monodromy-z3}
\end{figure}

\section{Formuła Picarda--Lefschetza i jej zakres}
\begin{twierdzenie}[Picard--Lefschetz dla rodzin krzywych]
\label{thm:picard-lefschetz-xiv}
Niech $\pi:\mathcal C\to\Delta$ będzie właściwym
holomorficznym odwzorowaniem, którego włókna nad
$\Delta^*$ są gładkimi zwartymi spójnymi
krzywymi zespolonymi, a włókno nad $0$
ma dokładnie jeden zwyczajny węzeł i poza nim
jest gładkie. Zakładamy, że całkowita przestrzeń
$\mathcal C$ jest gładka w otoczeniu węzła
i że w holomorficznych współrzędnych w tym otoczeniu
odwzorowanie $\pi$ ma lokalną postać
$\pi(z,w)=zw$ (przy współrzędnej $t$ na $\Delta$).
Niech $\delta\in H_1(C_t;\Z)$ będzie
zorientowanym cyklem zanikającym.
Po ustaleniu orientacji pętli i konwencji
sparowania przecięć monodromia na $H_1(C_t;\Z)$
ma postać transwekcji
\[
T(\gamma)=\gamma+\langle\gamma,\delta\rangle\,\delta.
\]
Znak przy ostatnim składniku zależy od
przyjętej konwencji orientacji.
\end{twierdzenie}
\begin{proof}[Status i rachunek lokalnego modelu]
Pełna identyfikacja monodromii globalnej z
skręceniem Dehna wokół $\delta$ jest twierdzeniem
Picarda--Lefschetza. Po jego przyjęciu wzór na
homologię wynika z policzenia wpływu skręcenia
na krzywą $\gamma$: każde poprzeczne przecięcie
z $\delta$ dokłada jeden obieg wokół $\delta$,
a znaki tych obiegów sumują się do
$\langle\gamma,\delta\rangle$.
Nie sprowadzamy tego globalnego rachunku do
stwierdzenia o monodromii lokalnego pierścienia
$\{zw=t\}$, gdzie działanie na $H_1$ jest identyczne.
Założenie o lokalnym modelu $\pi(z,w)=zw$ jest istotne:
po podstawieniu $t=s^m$ otrzymujemy rodzinę
$zw=s^m$ i $m$-tą potęgę monodromii, więc bez
tej hipotezy wzór na pojedynczą transwekcję
nie musi być poprawny.
\end{proof}
\begin{przyklad}[Macierz monodromii]
Załóżmy, że $C_t$ ma rodzaj $1$ oraz
$H_1(C_t;\Z)=\Z a\oplus\Z b$, przy czym
$\delta=a$ i $\langle b,a\rangle=1$.
Wówczas $T(a)=a$, $T(b)=b+a$.
W bazie $(a,b)$ macierz monodromii jest
$\bigl(\begin{smallmatrix}1&1\\0&1\end{smallmatrix}\bigr)$.
Jej nilpotentna część $N=T-I$ spełnia
$N\ne0$ oraz $N^2=0$; ilustruje to różnicę
między monodromią unipotentną a identyczną.
\end{przyklad}

\section{Powrót do snopów perwersyjnych i ograniczenia teorii}
\begin{twierdzenie}[Dokładność perwersyjna cykli]
Niech $X$ będzie zespoloną odmianą algebraiczną,
$f:X\to\C$ morfizmem algebraicznym,
$X_0=f^{-1}(0)$, a $P\in\operatorname{Perv}(X;\Q)$.
W konwencji Definicji~\ref{def:nearby-vanishing-xiv}
funktory $R\psi_f[-1]$ i $R\phi_f[-1]$
z $D_c^b(X;\Q)$ do $D_c^b(X_0;\Q)$
są $t$-dokładne względem środkowych
perwersyjnych struktur $t$; w szczególności
ich wartości na $P$ należą do
$\operatorname{Perv}(X_0;\Q)$.
\end{twierdzenie}
\begin{proof}[Status]
Wynik wymaga pełnej teorii funktorów cykli
pobliskich i zanikających oraz sklejenia
perwersyjnych struktur $t$. Przyjmujemy go
jako twierdzenie; powyższe rachunki dla
$A_1$ same w sobie nie dają jego ogólnego dowodu.
\end{proof}
\begin{uwaga}[Zastosowania i dalsza granica]
Cykle zanikające rejestrują topologię zmieniającą
się podczas degeneracji; w szczególności w
wymiarze zespolonym $n$ dla izolowanych punktów
krytycznych są skupione w wymiarze środkowym
(co wynika z twierdzenia Milnora o bukiecie sfer).
Pełny formalizm mieszanych modułów Hodge'a,
teoremat o quasi-unipotentności monodromii
oraz dowody twierdzenia o rozkładzie przez
czystość wag pozostają poza zakresem tego rozdziału.
Użycie słów ``twierdzenie'' i ``dowód'' w
poprzednich rozdziałach nie usuwa tych zależności.
\end{uwaga}

\section{Kontrola zależności i przygotowanie do geometrii algebraicznej}
W tym rozdziale bez odwoływania się do głębokich twierdzeń
dowiedliśmy tożsamości
$\dd^2=0$ dla stożka, jego długiego ciągu,
równoważności ``quasi-izomorfizm $\Leftrightarrow$
stożek acykliczny'', rachunku filtracji
czystej struktury Hodge'a oraz typu homotopijnego
i monodromii osobliwości kwadratowej. Pozostałe
twierdzenia oznaczyliśmy jako użyte narzędzia.
Przy scalaniu z geometrią algebraiczną
należy współdzielić definicje snopów, hiper-
kohomologii i funktorów pochodnych, a nie
powielać je w obu skryptach.

